% A tiny blueprint that exercises the reconciliation taxonomy.
% Pair it with axioms.txt (a saved axiom-report) -- no Lean build needed.
%
% Two DISTINCT things a green blueprint can hide, kept apart by the reconciler:
%   - a VALIDITY gap (kernel-refutes-claim): \leanok but #print axioms shows sorryAx
%       -> lem:helper, and thm:main which really inherits the hole.
%   - a COVERAGE gap (blueprint-incomplete): the Lean proof is genuinely clean, but the
%       blueprint \uses an unformalized concept -> thm:bridge. This is NOT a validity gap;
%       the theorem stays machine-checked. It is a documentation/faithfulness signal.

\begin{definition}[Base]\label{def:base}
  \lean{Demo.base}\leanok
  A fully formalized base definition.
\end{definition}

\begin{lemma}[Helper marked done but still open]\label{lem:helper}
  \lean{Demo.helper}\leanok
  \uses{def:base}
  The author marked this formalized, but the Lean proof still has a hole.
\end{lemma}

\begin{theorem}[Main result resting on the helper]\label{thm:main}
  \lean{Demo.main}\leanok
  \uses{lem:helper}
  Green in the blueprint, but its kernel axioms inherit the helper's hole.
\end{theorem}

\begin{theorem}[Bridge whose Lean proof is clean]\label{thm:bridge}
  \lean{Demo.bridge}\leanok
  \uses{thm:informal-step}
  The Lean proof is genuinely clean, yet the blueprint declares it rests on an
  unformalized step.
\end{theorem}

\begin{lemma}[Informal step, not formalized]\label{thm:informal-step}
  A paper-level statement that the bridge declares it uses, with no Lean
  declaration of its own.
\end{lemma}
